\begin{abstract}
What information must a change of representation preserve for a formal
proof to continue? This review studies that question in six developments
from the OpenAI Lean collection: common-base counting, graph switch
chains, the symmetric Mahler inequality, the complete Crouzeix bound,
extended trace cones, and marked character-variety constructions.
We follow the declarations that transport probability laws, computational
bounds, analytic estimates, and coordinate identities to their downstream
uses. The cases expose a division of proof work: mathematical interfaces
establish the hypotheses under which local arithmetic, algebraic, and
finite checks suffice. Elementary examples explain the information these
interfaces must retain. An exploratory survey compares 800 uniformly
sampled OpenAI modules with 8,070 mathematical modules in their pinned
Mathlib dependency. With equal weight per proof-bearing module and joint
standardization by broad subject and proof-length band, arithmetic and
algebra tactics occur in a larger share of extracted named proof bodies
by 22.8 percentage points; the differences are 10.4 points for tactic-mode
proofs and 5.8 for local claims or calculation chains. Comparing these
estimates with subject-only controls separates persistent arithmetic
emphasis from differences associated with proof length. Large numerical
certificate batches also make pooled body frequencies differ sharply
from module averages. Together, source tracing and corpus analysis
identify how established Lean methods are organized in this collection
and give readers a practical way to inspect the mathematical obligations
behind a formal theorem.
\end{abstract}
